\subsection{Polynomial Equations in Factored Form} 
\begin{fact}[Zero Property of Multiplication]
For any real numbers $a$ and $b$, if $ab=0$ then $a=0$ or $b=0$
\end{fact}
We can use this fact to solve polynomial equations.
\begin{Example}
Use the Zero Property of Multiplication to solve the polynomial equation
$$(2x-1)(x+3)(x-5)=0$$
Since the \makeblank[1.5in] is \makeblanksmall, one of the \makeblank[1.5in] is \makeblanksmall.
\begin{lpic}{}{3}
\end{lpic}
We can check that these really are solutions to the original equation:
\begin{lpic}{}{3}
\end{lpic}
\end{Example}
\begin{Note}
We will solve \emph{many} equations of the form $ax+b=0$ as we continue to work with polynomials! The solution is always $x=-\frac{b}{a}$.
\end{Note}